% from ega3-1-fr.tex
\oldpage[0\textsubscript{III}]{13}
$V\cap\complement{U}$, où $U$ est réunion d'un certain nombre des $U_i$ et
$V$ intersection d'un certain nombre des $V_i$~; mais on a remarqué plus
haut que les réunions et intersections finies d'ouverts rétrocompacts dans
$X$ sont des ouverts rétrocompacts dans $X$, d'où la conclusion.

\begin{corollary}[9.1.4]
\label{0.9.1.4-fr}
Toute partie constructible de $X$ est rétrocompacte dans $X$.
\end{corollary}

Il suffit de montrer que si $U$ et $V$ sont ouverts rétrocompacts dans $X$,
$U\cap\complement{V}$ est rétrocompact dans $X$~; or, si $W$ est ouvert
quasi-compact dans $X$, $W\cap U\cap\complement{V}$ est fermé dans l'espace
quasi-compact $W\cap U$, donc est quasi-compact.

En particulier~:

\begin{corollary}[9.1.5]
\label{0.9.1.5-fr}
Pour qu'une partie ouverte $U$ de $X$ soit constructible, il faut et il
suffit qu'elle soit rétrocompacte dans $X$. Pour qu'une partie fermée $F$ de
$X$ soit constructible, il faut et il suffit que l'ouvert $\complement{F}$
soit rétrocompact.
\end{corollary}

\begin{env}[9.1.6]
\label{0.9.1.6-fr}
Un cas important est celui où toute partie ouverte quasi-compacte de $X$ est
rétrocompacte, autrement dit, où l'intersection de deux parties ouvertes
quasi-compactes de $X$ est quasi-compacte
(cf. \hyperref[I.5.5.6-fr]{I, 5.5.6}). Lorsque $X$ lui-même est
quasi-compact, cela signifie que les parties ouvertes rétrocompactes dans
$X$ sont identiques aux parties ouvertes quasi-compactes de $X$, et les
parties constructibles de $X$ aux réunions finies d'ensembles de la forme
$U\cap\complement{V}$, où $U$ et $V$ sont ouverts quasi-compacts.
\end{env}

\begin{corollary}[9.1.7]
\label{0.9.1.7-fr}
Pour qu'une partie d'un espace noethérien $X$ soit constructible, il faut et
il suffit qu'elle soit réunion finie de parties localement fermées de $X$.
\end{corollary}

\begin{proposition}[9.1.8]
\label{0.9.1.8-fr}
Soient $X$ un espace topologique, $U$ une partie ouverte de $X$.
\begin{enumerate}
  \item[(i)] Si $T$ est une partie constructible de $X$, $T\cap U$ est une
    partie constructible de $U$.
  \item[(ii)] Supposons en outre $U$ rétrocompact dans $X$. Pour qu'une
    partie $Z$ de $U$ soit constructible dans $X$, il faut et il suffit
    qu'elle soit constructible dans $U$.
\end{enumerate}
\end{proposition}

\begin{enumerate}
  \item[(i)] Utilisant (\hyperref[0.9.1.3-fr]{9.1.3}), on est ramené à
    montrer que si $T$ est ouvert rétrocompact dans $X$, $T\cap U$ est ouvert
    rétrocompact dans $U$, autrement dit, pour tout ouvert quasi-compact
    $W\subset U$, $T\cap U\cap W=T\cap W$ est quasi-compact, ce qui résulte
    aussitôt de l'hypothèse.
  \item[(ii)] La condition étant nécessaire en vertu de (i), il reste à
    démontrer qu'elle est suffisante. Compte tenu de
    (\hyperref[0.9.1.3-fr]{9.1.3}), il suffit de considérer le cas où $Z$ est
    ouvert rétrocompact \emph{dans $U$}, car il s'ensuivra alors que
    $U\setminus Z$ est constructible dans $X$, et si $Z$, $Z'$ sont deux
    ouverts rétrocompacts \emph{dans $U$},
    $Z\cap(U\setminus Z')$ sera bien constructible dans $X$. Or, si $W$ est
    ouvert quasi-compact dans $X$ et $Z$ ouvert rétrocompact dans $U$, on a
    $Z\cap W=Z\cap(W\cap U)$ et par hypothèse $W\cap U$ est ouvert
    quasi-compact dans $U$~; donc $W\cap Z$ est bien quasi-compact, et par
    suite $Z$ est ouvert rétrocompact dans $X$, et \emph{a fortiori}
    constructible dans $X$.
\end{enumerate}

\begin{corollary}[9.1.9]
\label{0.9.1.9-fr}
Soient $X$ un espace topologique, $(U_i)_{i\in I}$ un recouvrement fini de
$X$ formé d'ensembles ouverts rétrocompacts dans $X$. Pour qu'une partie $Z$
de $X$ soit constructible dans $X$, il faut et il suffit que pour tout
$i\in I$, $Z\cap U_i$ soit constructible dans $U_i$.
\end{corollary}

\begin{env}[9.1.10]
\label{0.9.1.10-fr}
Supposons en particulier que $X$ soit quasi-compact et que tout point
